% !TEX root = ../1Main.tex
% Grundlagen - Template/Beispielkapitel

\chapter{Grundlagen}
\label{chap:grundlagen}

% ===================================================================
% ANLEITUNG: Definieren Sie hier die theoretischen Grundlagen.
% Dieses Kapitel zeigt Beispiele für verschiedene Formatierungen.
% ===================================================================

In diesem Kapitel werden die theoretischen Grundlagen und wichtige Begriffe eingeführt, die zum Verständnis der Arbeit notwendig sind.

\section{Begriffsdefinitionen}
\label{sec:begriffe}

Wichtige Begriffe sollten klar definiert werden. Sie können dafür die Definition-Umgebung nutzen:

\begin{definition}[Wichtiger Begriff]
    \label{def:beispiel}
    Ein \textbf{wichtiger Begriff} beschreibt \dots
\end{definition}

Später können Sie auf die Definition verweisen: Wie in \autoref{def:beispiel} definiert, beschreibt \dots

\subsection{Weitere Begriffe}

\begin{definition}[Zweiter Begriff]
    \label{def:zweiter-begriff}
    Diese Definition zeigt einen weiteren wichtigen Begriff.
\end{definition}

\section{Theoretischer Hintergrund}
\label{sec:theorie}

Hier beschreiben Sie die relevanten Theorien und Konzepte. Literaturverweise sind wichtig für wissenschaftliche Arbeiten.

\subsection{Grundlegende Konzepte}

Das Konzept X wurde erstmals in der Literatur beschrieben und später erweitert. Die wichtigsten Eigenschaften sind:
\begin{itemize}
    \item Eigenschaft 1: Beschreibung
    \item Eigenschaft 2: Beschreibung
    \item Eigenschaft 3: Beschreibung
\end{itemize}

\subsection{Verwandte Arbeiten}

Verwandte Arbeiten in diesem Bereich umfassen verschiedene Ansätze. Eine tabellarische Übersicht kann hilfreich sein:

\begin{table}[htbp]
    \centering
    \caption{Vergleich existierender Ansätze}
    \label{tab:vergleich-ansaetze}
    \begin{tabular}{lccc}
        \toprule
        \textbf{Ansatz} & \textbf{Vorteil} & \textbf{Nachteil} & \textbf{Jahr} \\
        \midrule
        Ansatz A & Schnell & Ungenau & 2020 \\
        Ansatz B & Genau & Langsam & 2021 \\
        Ansatz C & Ausgewogen & Komplex & 2022 \\
        \bottomrule
    \end{tabular}
\end{table}

Wie in \autoref{tab:vergleich-ansaetze} dargestellt, haben alle Ansätze spezifische Vor- und Nachteile.

\section{Mathematische Grundlagen}
\label{sec:mathematik}

Falls Ihre Arbeit mathematische Formeln enthält, können Sie diese inline darstellen $f(x) = x^2 + 2x + 1$ oder als separate Gleichung:

\begin{equation}
    \label{eq:beispiel}
    f(x) = \sum_{i=0}^{n} a_i x^i
\end{equation}

Verweise auf Gleichungen erfolgen mit \autoref{eq:beispiel}.

Für mehrzeilige Gleichungen nutzen Sie die \texttt{align}-Umgebung:

\begin{align}
    y &= mx + b \\
    \frac{dy}{dx} &= m
\end{align}

\subsection{Wichtige Theoreme}

\begin{satz}[Beispielsatz]
    \label{satz:beispiel}
    Für alle $x \in \mathbb{R}$ gilt: $x^2 \geq 0$
\end{satz}

\begin{lemma}
    \label{lemma:beispiel}
    Dies ist ein Beispiel für ein Lemma, das als Hilfssatz dient.
\end{lemma}

\section{Technologische Grundlagen}
\label{sec:technologie}

Beschreiben Sie hier relevante Technologien, Frameworks oder Werkzeuge.

\subsection{Technologie A}

Technologie A bietet folgende Features:
\begin{enumerate}
    \item Feature 1
    \item Feature 2
    \item Feature 3
\end{enumerate}

\subsection{Technologie B}

Technologie B wird verwendet, weil \dots

% Beispiel für Code in den Grundlagen (kurzer Ausschnitt)
Code-Beispiele können inline mit \texttt{variable = 42} oder als Block formatiert werden.

\section{Zusammenfassung}
\label{sec:grundlagen-zusammenfassung}

Dieses Kapitel hat folgende Aspekte behandelt:
\begin{itemize}
    \item Wichtige Begriffsdefinitionen (\autoref{sec:begriffe})
    \item Theoretischer Hintergrund (\autoref{sec:theorie})
    \item Mathematische Grundlagen (\autoref{sec:mathematik})
    \item Technologische Grundlagen (\autoref{sec:technologie})
\end{itemize}

Diese Grundlagen bilden die Basis für die weiteren Kapitel dieser Arbeit.

% ===================================================================
% HINWEIS: Passen Sie die Inhalte an Ihr Thema an!
% ===================================================================